\section{\seclogicconn}
\label{sec_logic_connectives}

In this section, we introduce the logical connectives for ``and'' and ``or'' along with their proof formats.  We discuss logical equivalence and negation and use truth tables to define connectives and verify equivalences.

Let's warm up our logical thinking by considering some true/false questions.

\begin{act}[True or false?]
\label{act_TorF}
    Decide if each statement is true or false and discuss how to justify your answer.

        \begin{actpart}
            \item \label{part_digit_sum23} There are three distinct digits whose sum is 23. %True THERE EXIST
            \item \label{part_digit_sum25} There are three distinct digits whose sum is 25. %False THERE EXIST
            \item \label{part_pos_or_neg} Every integer is positive or negative. %False % FOR ALL OR
            \item \label{part_div3_and_div5} Some integers are divisible by three and five. %True EX AND
            \item \label{part_div6_but_ndiv3} Some integers are divisible by six but not by three. %False  EX But
            \item \label{part_graph_order5} A graph with five vertices is the cycle $C_5$ or the complete graph $K_5$. % False OR
            \item Every graph is a tree, which means that it is connected and acyclic. % False AND
            \item Every bit string of length five begins or ends with \bs{0}. % False % FOR ALL OR
            \item \label{part_bs5_three0_or_three1} Every bit string of length five has at least three \bs{0}s or at least three \bs{1}s. %True % FOR ALL OR EXISTS
        \end{actpart}
\end{act}

%%%%%%%%%%%%%%%%%%%%%%%%%%%%%%%%%%%%%%%%%%
\newcommand{\subandor}{And and Or}
\subsection{\subandor}
\label{sub_and_or}

In \Cref{act_TorF}, there are many words that indicate the logical structure of a statement.  Two important such words are ``and'' and ``or'', each of which has a specific technical meaning in mathematics and computer science.

\begin{defn}[And, or, and xor]
\label{defn_and_or_xor}
~
    \begin{apart}
        \item A \vocab{statement} is a sentence that is either true or false.  For example,
            \begin{center}
                $P=$ the integer six has exactly four positive divisors.
            \end{center}
        is a statement that happens to be true because the positive divisors of six are 1, 2, 3, and 6.  The statement
            \begin{center}
                $Q=$ the cycle $C_5$ has a vertex of degree three.   
            \end{center} 
        is false because all the vertices of $C_5$ have degree two.
        
        \item For statements $P$ and $Q$, the statement ``$P$ \vocab{and} $Q$'', denoted $P \land Q$, is true if both $P$ and $Q$ are true and false otherwise.  For example, the statement  
            \begin{center}
                The integer seven is positive and odd.    
            \end{center}
        is true because seven is positive and seven is odd. The statement 
            \begin{center}
                The integer six is positive and odd.   
            \end{center} 
        is false, even though six is positive because six is not odd.
        
        \item For statements $P$ and $Q$, the statement ``$P$ \vocab{or} $Q$'', denoted $P \lor Q$, is true if at least one of $P$ and $Q$ are true and false otherwise.  For example, the statement
            \begin{center}
                The integer six is positive or odd.
            \end{center}
        is true, because six is positive.  The statement 
            \begin{center}
                A graph with five vertices is either the cycle $C_5$ or the complete graph $K_5$.
            \end{center}
        from \Cref{act_TorF} (\ref{part_graph_order5}) is false because, for example, the graph $P_5$ has five vertices but it is not a cycle and it is not complete.  
        
        \item Note that $P \lor Q$ is true if both $P$ and $Q$ are true.  For example, the statement 
            \begin{center}
                The integer seven is positive or odd.
            \end{center} 
        is true.  This use of ``or'' is the \vocab{inclusive or} because it includes the possibility of both statements being true.
        
        \item When we want to exclude the possibility of both statements being true, we use the \vocab{exclusive or} instead. For statements $P$ and $Q$, the statement ``$P$ \vocab{xor} $Q$'', denoted $P \oplus Q$, is true if one of $P$ and $Q$ is true and the other is false and false otherwise.  For example, when I told my kids ``you can have ice cream or brownies for dessert'' they suggested brownies with ice cream (the inclusive or), but I probably meant one or the other but not both (the exclusive or).  As another example, in \Cref{exam_every_integer_even_or_odd} 
        we showed that 
            \begin{center}
                Every integer is either even or odd, but never both.
            \end{center}
        which we could write more concisely as
            \begin{center}
                Every integer is even xor odd.
            \end{center} 

        \item One way to summarize when a compound statement is true or false is to use a \vocab{truth table} which includes a column for each component of the compound statement (such as $P$, $Q$) and for the full compound statement (such as $P \land Q$) and  a row for each possible combination of true or false for the component statements.  For example, if our components are $P$ and $Q$, then the truth table will have four rows corresponding to the four possibilities:
            \begin{itemize}
                \item $P$ is true and $Q$ is true;
                \item $P$ is true and $Q$ is false;
                \item $P$ is false and $Q$ is true; and
                \item $P$ is false and $Q$ is false.
            \end{itemize}
        In this textbook we abbreviate true as \str{T} and false as \str{F}.  Some computer science textbooks use \str{1} for true and \str{0} for false instead. 
    \end{apart}   
\end{defn}

We show the truth tables for the compound statements $P \land Q$, $P \lor Q$, and $P \oplus Q$ in our next example.

\begin{exam}[Truth tables for and, or, and xor]
\label{exam_truthtable_and_or_xor}
~  
        \begin{apart}
            \item Use a truth table to illustrate the meaning of the logical connective $\land$ (and).
            \begin{soln}               
                \Cref{table_truthtable_and} shows the truth table for $P \land Q$.  Since the only time $P \land Q$ is true is when both $P$ and $Q$ are true, the only \str{T} in the $P \land Q$ column is in the first row.  The other entries in the $P \land Q$ column are all \str{F}.
                
                \begin{table} [hbtp]
                    \centering
                    \begin{tabular}{cc|c}
                        $P$ & $Q$ & $P \land Q$ \\ \hline
                        \str{T} & \str{T} & \str{T}\\
                        \str{T} & \str{F} & \str{F}\\
                        \str{F} & \str{T} & \str{F}\\
                        \str{F} & \str{F} & \str{F}\\
                    \end{tabular}
                    \caption{Truth table for the logical connective \vocab{and}.}
                    \label{table_truthtable_and}
                \end{table}
                
            \end{soln}
        \item Use a truth table to illustrate the meaning of the logical connective $\lor$ (or). 
            \begin{soln}
                \Cref{table_truthtable_or} shows the truth table for $P \lor Q$. Since $P \lor Q$ is true is whenever $P$ is true, the first two entries in the $P \land Q$ column are \str{T}.  Since $P \lor Q$ is true whenever $Q$ is true, the third entry in the $P \land Q$ column is also \str{T}. The final entry in the $P \land Q$ column is \str{F} because neither $P$ nor $Q$ are true.
                
                    \begin{table} [hbtp]
                    \centering
                    \begin{tabular}{cc|c}
                        $P$ & $Q$ & $P \lor Q$ \\ \hline
                        \str{T} & \str{T} & \str{T}\\
                        \str{T} & \str{F} & \str{T}\\
                        \str{F} & \str{T} & \str{T}\\
                        \str{F} & \str{F} & \str{F}\\
                    \end{tabular}
                    \caption{Truth table for the logical connective \vocab{or}.}
                    \label{table_truthtable_or}
                \end{table}
                
            \end{soln}
        \item Use a truth table to illustrate the meaning of the logical connective $\oplus$ (xor).
            \begin{soln}
                \Cref{table_truthtable_xor} shows the truth table for $P \oplus Q$. Since $P \oplus Q$ is true exactly when one of $P$ or $Q$ is true, the $P \oplus Q$ column has \str{T} in the second and third rows only.
                    \begin{table}[hbtp]
                    \centering
                    \begin{tabular}{cc|c}
                        $P$ & $Q$ & $P \oplus Q$ \\ \hline
                        \str{T} & \str{T} & \str{F}\\
                        \str{T} & \str{F} & \str{T}\\
                        \str{F} & \str{T} & \str{T}\\
                        \str{F} & \str{F} & \str{F}\\
                    \end{tabular}
                    \caption{Truth table for the logical connective \vocab{xor}.}
                    \label{table_truthtable_xor}
                \end{table}
                
            \end{soln}
        \end{apart} 
\end{exam}

It is your turn to create a truth table.

\begin{act}[Truth tables with and and or]
\label{act_truthtable_and_or} 
    \Cref{table_truthtable_and_or} shows the beginning of the truth table for the compound statement $(P \lor Q) \land R$.
            \begin{table} [hbtp]
                \centering
                \begin{tabular}{ccc|cc}
                    $P$ & $Q$ & $R$ & $P \lor Q$ & $(P \lor Q) \land R$\\ \hline
                    \str{T} & \str{T} & \str{T} & \str{T} & \str{T} \\
                    \str{T} & \str{T} & \str{F} & \str{T} & \str{F} \\
                    \str{T} & \str{F} & \str{T} & \str{T} & \str{T} \\
                    \str{T} & \str{F} & \str{F} & \str{T} & \str{F} \\
                    \str{F} & \str{T} & \str{T} & & \\
                    \str{F} & \str{T} & \str{F} & & \\
                    &&&& \\
                    &&&& \\
                \end{tabular}
                \caption{Part of the truth table for $(P \lor Q) \land R$}
                \label{table_truthtable_and_or}
            \end{table}
    \begin{actpart}
        \item According to the truth table in \Cref{table_truthtable_and_or}, when $P$ is true, $Q$ is false, and $R$ is true, which row are we in (do not count the ``header row'' with $P$, $Q$, etc.) and is the statement $(P \lor Q) \land R$ true or false?
        \item We are missing the last couple of rows of \Cref{table_truthtable_and_or}.  What are the remaining true/false options?  Fill in the $P$, $Q$, and $R$ columns for the missing rows. 
        \item In \Cref{table_truthtable_and_or} we included an intermediate column for $P \lor Q$.  Convince yourself that we correctly started that column and complete it.  Notice that we ignore the $R$ column.
        \item \label{part_finish_truthtable} Convince yourself that we correctly started the final column for $(P \lor Q) \land R$ in \Cref{table_truthtable_and_or} and complete it.  Notice that we need to look at only the column for $R$ and the intermediate column for $P \lor Q$.
        \item Construct a new truth table for the compound statement $(P \land R) \lor (Q \land R)$. Include intermediate columns for $P \land R$ and for $Q \land R$.  What do you notice?  Hint: Compare to \Cref{table_truthtable_and_or}.  

        \item If we wanted to construct a truth table for a compound statement involving four components ($P$, $Q$, $R$, and $S$), how many rows would it need.  Explain how you counted.  Hint: Steps.
    \end{actpart}
\end{act}

Make sure that you have completed \Cref{act_truthtable_and_or} because we are going to reveal some of the answers.

\begin{exam}[Distributive rule]
\label{exam_distributive_equiv}
    Use a truth table to verify the logical equivalence $$ (P \lor Q) \land R \equiv (P \land R) \lor (Q \land R).$$
        \begin{soln}
            The truth table is in \Cref{table_truthtable_distrib}.  Notice that the columns for $(P \lor Q) \land R$ and for $(P \land R) \lor (Q \land R)$ have the same sequence of \str{T}s and \str{F}s.
            
            \begin{table} [hbtp]
                \centering
                \begin{tabular}{ccc|cc|ccc}
                    $P$ & $Q$ & $R$ & $P \lor Q$ & $(P \lor Q) \land R$ & $P \land R$ & $Q \land R$ & $(P \land R) \lor (Q \land R)$
                    \\ \hline
                    \str{T} & \str{T} & \str{T} & \str{T} & \str{T} & \str{T} & \str{T} & \str{T} 
                    \\
                    \str{T} & \str{T} & \str{F} & \str{T} & \str{F} & \str{F} & \str{F} & \str{F} 
                    \\
                    \str{T} & \str{F} & \str{T} & \str{T} & \str{T} & \str{T} & \str{F} & \str{T} 
                    \\
                    \str{T} & \str{F} & \str{F} & \str{T} & \str{F} & \str{F} & \str{F} & \str{F} 
                    \\
                    \str{F} & \str{T} & \str{T} & \str{T} & \str{T} & \str{F} & \str{T} & \str{T} 
                    \\
                    \str{F} & \str{T} & \str{F} & \str{T} & \str{F} & \str{F} & \str{F} & \str{F} 
                    \\
                    \str{F} & \str{F} & \str{T} & \str{F} & \str{F} & \str{F} & \str{F} & \str{F} 
                    \\
                    \str{F} & \str{F} & \str{F} & \str{F} & \str{F} & \str{F} & \str{F} & \str{F} 
                \end{tabular}
                \caption{A truth table verifying the distributive equivalence $ (P \lor Q) \land R \equiv (P \land R) \lor (Q \land R)$.}
                \label{table_truthtable_distrib}
            \end{table}
        \end{soln}
\end{exam}
%%%%%%%%%%%%%%%%%%%%%%%%%%%%%%%%%%%%%%%%%%
\newcommand{\subpffandor}{Proof Formats: $\land$ and $\lor$ by Cases}
\subsection{\subpffandor}
\label{sub_pff_and_or}

How can we prove statements involving and and or? To prove $P \land Q$, we need to prove two things: $P$ and then $Q$.  We can summarize this idea in a proof format.

\begin{pff}[Direct proof of $\land$]
\label{pff_and}
    Prove: $P \land Q$
        \begin{proof}            
            First, (explain why $P$ is true.)  Therefore, $P$.

            Next, (explain why $Q$ is true.) Therefore, $Q$.
        \end{proof}
\end{pff}

The format for a proof of $P \lor Q$ is less direct. An option is to use a modified version of a proof by cases. Later in this textbook, we introduce two other options: proof by contradiction in \Cref{sub_pff_contra} and proof using the equivalent conditional in \Cref{sub_pff_if_then}.

In \Cref{sub_pff_cases} we introduced proof by cases.  In each case of the format in \Cref{pff_cases}, we proved the same statement.  We can modify that proof format to prove $P \lor Q$ by proving $P$ in one case and $Q$ in the other case.

\begin{pff}[Proof of $\lor$ by cases]
\label{pff_or_by_cases}
    Prove: $P \lor Q$
        \begin{proof}   
            (Explain why there are two cases.)  We consider two cases.

            Case 1: Assume \ldots. (Explain why $P$ is true in this case.)  Therefore, $P$.

            Case 2: Assume \ldots.  (Explain why $Q$ is true in this case.) Therefore, $Q$.

            In either case, we have $P \lor Q$.
        \end{proof}
\end{pff}

We present this format with two cases, since that is often all we need.  This format can be modified to have more than two cases where in some cases $P$ is true and in other cases $Q$ is true. Officially, this is a new proof format.

\begin{pff} [Proof of $\lor$ using cases, generalized]
\label{pff_or_by_cases_generalized} 
Prove: $P_1 \lor P_2 \lor \cdots \lor P_k$
    \begin{proof}
        (Explain why there are $n$ cases.) We consider $n$ cases.
  
        Case 1:  Assume \ldots (Explain why one of $P_1, P_2, \ldots P_k$ is true in this case.) 
        
        Case 2:  Assume \ldots (Explain why one of $P_1, P_2, \ldots P_k$ is true in this case.) 
        
        \ldots

        Case $n$: Assume \ldots (Explain why one of $P_1, P_2, \ldots P_k$ is true in this case.) 
        
        In (either case/all cases) we have $P_1 \lor P_2 \lor \cdots \lor P_k$. 
\end{proof}
\end{pff}

Try working with these formats.

\begin{act}[Proving $\land$ and $\lor$ by using cases]
\label{act_prove_and_orusingcases}
~
    \begin{actpart}
        \item \label{part_pff_and_example} Copy the following example of a proof of an ``and'' statement.

        Prove that 30 is even and $5 \mid 30$.

        \begin{proof} First, $30 = 2(15)=2k$ when $k=15$ and so 30 is even. 

        Next, $30 = 5(6) = 5k$ when $k=6$ and so $5 \mid 30$.\end{proof}
        
        \item Edit the proof in (\ref{part_pff_and_example}) to prove that 33 is odd and $3 \mid 33$.

        \item \label{part_pff_or_using_cases_example} Copy the following example of a proof of an ``or'' statement using cases.

        Prove that $n^2 \fn{ mod } 4 = 0$ or $n^2 \fn{ mod } 4 = 1$ for any integer $n$. 
        
        Yes, this proof was \Cref{exer_proof_cases_square_mod4} in \Cref{sec_div_alg}.

        \begin{proof}
            Let $n$ be an integer.  Then $n$ is either even or odd, by \Cref{exam_every_integer_even_or_odd}.  We consider two cases.

            Case 1: Assume that $n$ is even.  By the definition of even, we can write $n = 2k$ for some integer $k$.  Then 
                $$n^2=(2k)^2=4k^2=4(k^2) = 4q=4q+0$$
            where $q=k^2$ is an integer.  Therefore, $n^2 \fn{ mod } 4 = 0.$

            Case 2: Assume that $n$ is odd.  By the definition of odd, we can write $n = 2k+1$ for some integer $k$.  Then 
                $$n^2=(2k+1)^2 = (2k+1)(2k+1)=4k^2+4k+1 = 4(k^2+k) + 1 = 4q+1$$
            where $q =k^2+k$ is an integer.  Therefore, $n^2 \fn{ mod } 4 = 1.$

            In either case, we have $n^2 \fn{ mod } 4 = 0$ or $n^2 \fn{ mod } 4 = 1$.
        \end{proof}

        \item \label{part_nsqplusn_mod4} Edit the proof in (\ref{part_pff_or_using_cases_example}) to prove that $n^2+2n \fn{ mod } 4 = 0$ or $n^2+2n \fn{ mod } 4 = 3$ for any integer $n$.  
    \end{actpart}
\end{act}

%%%%%%%%%%%%%%%%%%%%%%%%%%%%%%%%%%%%%%%%%%
\newcommand{\subequivneg}{Logical Equivalence and Negations}
\subsection{\subequivneg}
\label{sub_logical_equivalence_negations}

In \Cref{act_truthtable_and_or}, we saw that the statements $(P \lor Q) \land R$ and $(P \land R) \lor (Q \land R)$ have the same final column in their truth table.  That is, in each case (row of the truth table), they are both true or both false.  We have a name for the situation where two logical statements have the same truth values.

\begin{defn}[Logical equivalence]
\label{defn_logical_equiv}  
~
    \begin{apart}
        \item \label{part_equiv_distributive} Statements $S$ and $T$ are \vocab{logically equivalent}, denoted $S \equiv T$, if $S$ and $T$ always have the same truth value, either both statements are true or both statements are false.  For example, we saw in \Cref{act_truthtable_and_or} that 
            $$(P \lor Q) \land R \equiv (P \land R) \lor (Q \land R).$$
        \item Statements $S$ and $T$ are \vocab{not logically equivalent}, denoted $S \not \equiv T$ if sometimes one of the statements is true but the other is false.  For example,  
            $$(P \lor Q) \land R \not \equiv P \lor (Q \land R)$$
        as we check in \Cref{exer_not_assoc} because, for example, when $P$ is true and both $Q$ and $R$ are false, the left-hand side $(P \lor Q) \land R$ is false but the right-hand side $P \lor (Q \land R)$ is true.
    \end{apart}   
\end{defn}

Logically equivalent statements always have the same truth values.  There is also a name for when two statements always have the opposite truth value as well as names for logical statements that are always true or always false.

\begin{defn}[Negation, tautology, and contradiction]
\label{defn_negation_tautology_contradiction}
~
    \begin{apart}
        \item The statements $S$ and $T$ are \vocab{negations} if $S$ and $T$ always have the opposite truth value. That is, $S$ is true and $T$ is false or $S$ is false and $T$ is true.  The statement $T$ is the \vocab{negation} of $S$, denoted $\neg S$.  For example, if $n$ is an integer and if $S$ is the statement $n >2$, then $\neg S$ is the statement $n \le 2$.

        \item A statement $S$ is a \vocab{tautology} if it is always true.  For example, the statement 
            \begin{center}
                Either I will work on my homework Tuesday night or I will not work on my homework Tuesday night.   
            \end{center} 
        is a tautology.  

        \item A statement $S$ is a \vocab{contradiction} if it is always false.  For example, the statement 
            \begin{center}
                There is an integer $n$ with $n>2$ and $n^2<4$.    
            \end{center} 
        is a contradiction because for an integer $n$ with $n>2$ we know that $n^2 > 4$. 
    \end{apart}   
\end{defn}

Here are examples of truth tables involving negations.

\begin{exam}[Truth tables involving negations]
\label{exam_truthtable_negations}
~
    \begin{apart}
        \item Use a truth table to illustrate the meaning of the logical operation $\neg$ (not). 
            \begin{soln}
                \Cref{table_truthtable_not} shows the truth table for $\neg P$. 
            \end{soln}  
            
            \begin{table}[hbtp]
                \centering
                \begin{tabular}{c|c}
                        $P$ & $\neg P$  \\ \hline
                        \str{T} & \str{F} \\
                        \str{F} & \str{T} 
                \end{tabular}
                \caption{Truth table for the logical operation \vocab{not}.}
                \label{table_truthtable_not}
            \end{table}
            
        \item Construct a truth table for the compound statement $P \land \neg Q$.
            \begin{soln}
                \Cref{table_truthtable_andnot} shows the truth table for $P \land \neg Q$, using an intermediate column for $\neg Q$.  Notice that the only time $P \land \neg Q$ is true is when $P$ is true and $Q$ is false, which happens in the second row.
            \end{soln}
            
            \begin{table}[hbtp]
                \centering
                \begin{tabular}{cc|c|c}
                    $P$ & $Q$ & $\neg Q$ & $P \land \neg Q$ \\ \hline
                    \str{T} & \str{T} & \str{F} & \str{F}\\
                    \str{T} & \str{F} & \str{T} & \str{F}\\
                    \str{F} & \str{T} & \str{F} & \str{F}\\
                    \str{F} & \str{F} & \str{T}& \str{F}
                \end{tabular}
                \caption{Truth table for the compound statement $P \land \neg Q$.}
                \label{table_truthtable_andnot}
            \end{table}
        
        \item Use a truth table to show that $P \land \neg P$ is a contradiction.
            \begin{soln}
            \Cref{table_truthtable_contradiction} shows the truth table for $P \land \neg P$. Since the final column is all \str{F}, it follows that $P \land \neg P$ is a contradiction.
            \end{soln} 
                
            \begin{table}[hbtp]
                \centering
                \begin{tabular}{c|c|c}
                    $P$ & $\neg P$ & $P \land \neg P$ \\ \hline
                    \str{T} & \str{F} & \str{F}\\
                    \str{F} & \str{T} & \str{F}
                \end{tabular}
                \caption{Truth table for the contradiction $P \land \neg P$.}
                \label{table_truthtable_contradiction}
                \end{table}                   
    \end{apart} 
\end{exam}

By the way, when the statement following the logical connective ``and'' is negative, we use the connective ``but'' instead.

\begin{rem}[But not]
\label{rem_but_not}
   It is common in English to pronounce $P \land \neg Q$ as ``$P$ but not $Q$''.  The word ``but'' has the same meaning logically as ``and'' but in English it indicates that what follows is negative in some way.  For example, while we can say 
       \begin{center}
           I like chocolate ice cream and I don't want any now.    
       \end{center} 
    it is more common to say 
        \begin{center}
            I like chocolate ice cream but I don't want any now.
        \end{center}   
\end{rem}

Try your hand at working with negations.

\begin{act}[Negating and]
\label{act_not_and}
~
    \begin{actpart}
        \item In \Cref{defn_and_or_xor}, we noted that the statement  
            \begin{center}
                The integer seven is positive and odd.
            \end{center} 
        is true.  Consider the general statement 
            \begin{center}
                The integer $n$ is positive and odd.
            \end{center}  
            Give examples of integers $n$ where the statement is false, including some negative values of $n$.
        \item \label{part_positive_and_odd} For the statement 
            \begin{center}
                The integer $n$ is positive and odd.   
            \end{center} 
        to be false, what needs to be true for $n$?
        \item \label{part_truth_not_and} Construct a truth table for $\neg(P \land Q)$ and $\neg P \land \neg Q$.  
        \item Look at your truth table from (\ref{part_truth_not_and}).  Is $\neg(P \land Q) \equiv \neg P \land \neg Q$?  Explain.
        \item Is $\neg P \land \neg Q$ the negation of $P \land Q$? Explain.
        \item \label{part_conj_not_and} Make a conjecture about the negation of $P \land Q$. Your answer should agree with your answer to (\ref{part_positive_and_odd}).
        \item Verify your conjecture from (\ref{part_conj_not_and}) using a truth table.
    \end{actpart}
\end{act}

In \Cref{act_not_and}, we saw that the negation of $P \land Q$ is not $\neg P \land \neg Q$.  For $P \land Q$ to be false, we need either $P$ is false or $Q$ is false or both are false. That is, we need $\neg P \lor \neg Q$ to be true. Similarly, for $P \lor Q$ to be false, we need both $P$ is false and $Q$ is false.  That is, we need $\neg P \land \neg Q$ to be true.  We summarize these negations in a theorem.

\begin{thm}[De Morgan's Law]
\label{thm_demorgans}
For statements $P$ and $Q$, the negation of $P \land Q$ is $\neg P \lor \neg Q$ and the negation of $P \lor Q$ is $\neg P \land \neg Q$.  That is, we have the following logical equivalences known as \vocab{De Morgan's Laws}:
$$\neg(P \land Q) \equiv \neg P \lor \neg Q\text{ and }\neg(P \lor Q) \equiv \neg P \land \neg Q.$$  
\end{thm}

%%%%%%%%%%%%%%%%%%%%%%%%%%%%%%%%%%%%%%%%%%
\newcommand{\subpffcontra}{Proof Format: Contradiction}
\subsection{\subpffcontra}
\label{sub_pff_contra}

One method of proving that the statement $P$ is true is to show that the assumption that $P$ is false leads to a contradiction, such as a statement of the form $Q \land \neg Q$.  That is, in a proof by contradiction of a statement $P$, we assume that $P$ is false and deduce that both $Q$ and $\neg Q$ are true for some statement $Q$, which is impossible. That impossibility tells us that $P$ cannot be false, and so $P$ must be true.  

As we start writing a proof by contradiction, we usually do not know what the statement $Q$ will be.  Sometimes it helps to do a little scratch work listing the consequences of $P$ being false to see if we find contradictory statements $Q$ and $\neg Q$.

A \vocab{proof by contradiction} follows the following proof format.

\begin{pff}[Proof by contradiction]
\label{pff_contradiction}
    Prove: P
        \begin{proof}
            Suppose not. Then (state $\neg P$).

            First, (explain why $Q$ is true.)  Therefore, $Q$.

            On the other hand, (explain why $\neg Q$ is true.) Therefore, $\neg Q$.

            This is a contradiction. Therefore, $P$.
        \end{proof}
\end{pff}

Notice that we are proving two things ($Q$ and $\neg Q$), but instead of saying ``First, \ldots. Next, ldots.'' as we did in \Cref{pff_and}, we use the expression ``First, \ldots. On the other hand, ldots.'' The phrase ``on the other hand`` alerts the reader that we are changing direction towards the negation of what we just proved.  It is also acceptable to write ``On the one hand, \ldots. On the other hand, \ldots."

Try writing a proof using contradiction.

\begin{act}[Proof by contradiction: consecutive integers are coprime]
\label{act_proof_contradiction}
~
    \begin{actpart}
        \item \label{part_contradiction_consec_coprime} Copy the following example of a proof by contradiction.

        Prove $\fn{gcd}(n,n+1)=1$ for any integer $n$.\footnote{We proved this result in \Cref{exer_consec_coprime} using the Euclidean Algorithm.}   

        \begin{proof} Suppose not. Then there exists an integer $n$ such that $d=\fn{gcd}(n,n+1)\neq 1$. 
        
        First, by the definition of the greatest common divisor (\Cref{defn_gcd}), $d \mid n$ and $d \mid n+1$.  By the definition of divides (\Cref{defn_divides}) we can write $n = dk_1$ and $n+1 = dk_2$ for some integers $k_1$ and $k_2$.  Observe that $$(n+1) - n = dk_2 - dk_1 = d(k_2-k_1)=dq$$ where $q = k_2-k_1$ and so $d \mid (n+1)-n$.  But, of course, $(n+1)-n=1$ so really $d \mid 1$. Since the only positive divisor of 1 is 1, it follows that $d =1$.  

        On the other hand, we assumed that $d \neq 1$.  This is a contradiction. Thus, $\fn{gcd}(n,n+1)=1$.
        \end{proof}
     
         \item Edit the proof in (\ref{part_contradiction_consec_coprime}) to prove that $\fn{gcd}(n,n-1) = 1$ for any integer $n$. (Yes, we are again proving that consecutive integers are coprime.)
    \end{actpart}
\end{act}

There is a bit of controversy over the use of proof by contradiction.  Some mathematicians dislike contradiction proofs because they do not actually show why a statement is true, only why it cannot be false.  Other mathematicians like contradiction proofs because they settle an argument: if your reader thinks the statement is false, your proof shows them that something terrible goes wrong (a contradiction). We will let you decide which type of mathematician you are.

Because of this controversy, we often rewrite a proof by contradiction using some other format, if possible.  

\begin{exam}[Rewriting without contradiction]
\label{exam_avoid_contradiction_consec_coprime}
    Copy the proof by contradiction from \Cref{act_proof_contradiction} (\ref{part_contradiction_consec_coprime}) and edit to to give a proof without contradiction.  
    
    Prove that $\fn{gcd}(n,n+1)=1$ for any integer $n$.
    
        \begin{soln}
            We have crossed out the statements that should be deleted, and boldfaced the statements that should be added.

            \emph{Proof.} \st{Suppose not. Then there exists an integer $n$ such that $d=\fn{gcd}(n,n+1)\neq 1$} 
        
            \textbf{Let $n$ be an integer. Write $d=\fn{gcd}(n,n+1)$. Since 1 is a common divisor of any integer, $d$ is certainly positive.}
        
            First, by the definition of greatest common divisor (\Cref{defn_gcd}), $d \mid n$ and $d \mid n+1$.  By the definition of divides (\Cref{defn_divides}) we can write $n = dk_1$ and $n+1 = dk_2$ for some integers $k_1$ and $k_2$.  Observe that $$(n+1) - n = dk_2 - dk_1 = d(k_2-k_1)$$ and so $d \mid (n+1)-n$.  But, of course, $(n+1)-n=1$ so really $d \mid 1$. Since the only positive divisor of 1 is 1, it follows that $d =1$.  

            \st{On the other hand, we assumed that $d \neq 1$.  This is a contradiction. Thus $\fn{gcd}(n,n+1)=1$.} 
        \end{soln}    
\end{exam} 

%%%%%%%%%%%%%%%%%%%%%%%%%%%%%%%%%%%%%%%%%%
\subsection{Exercises}

%%%%%%%%%%%%%%%%%%%%%%%%

\subsubsection*{Exercises for \Cref{sub_and_or} \subandor}

\begin{exer}[\exerP] % standard truth tables]
\label{exer_standard_truthtables}
    Copy the standard truth tables for each expression. 
        \begin{exerpart}
            \item $P \land Q$
            \item $P \lor Q$
            \item $P \oplus Q$ 
        \end{exerpart}
\end{exer}

\begin{exer}[\exerP] % translating into statements using $\land$, $\lor$, and $\oplus$]
\label{exer_translate_into_and_or}
    For this exercise, $n$ is a specific (but unknown) integer,       
        \begin{center}
            $X=$ $n$ is a multiple of 6, and $R=$ $n$ is a multiple of 4.
        \end{center}  
    In each part of this exercise write the statement using $X$, $R$, and logic symbols and find an example of an integer $n$ for which the statement is true. 
        \begin{exerpart}
            \item The integer $n$ is a multiple 6 and a multiple of 4.
            \item The integer $n$ is a multiple of exactly one of 6 or 4.
        \end{exerpart}
\end{exer}

\begin{exer}[\exerU] % defns and or
    Recall that a graph is a \vocab{tree} if it is connected and acyclic. 
        \begin{exerpart}
            \item Give an example of a tree and explain what you have checked.
            \item Could a graph be connected but not a tree?  Give an example and explain in words.
            \item Could a graph be acyclic but not a tree?  Give an example and explain in words.
    \end{exerpart}
\end{exer}

\begin{exer}[\exerU]
\label{exer_not_assoc}  
~
    \begin{exerpart}
        \item \label{part_not_assoc} Construct a truth table for the compound statement $P \lor (Q \land R)$.  Include an intermediate column for $Q \land R$.
        \item Compare your answer from (\ref{part_not_assoc}) with the completed truth table from \Cref{act_truthtable_and_or} (\ref{part_finish_truthtable}). Is $P \lor (Q \land R) \equiv (P \lor Q) \land R$?  Explain.
    \end{exerpart}
\end{exer}

\begin{exer}[\exerU] % truth table for $P \land (Q \lor R)$] 
\label{exer_truthtable_distributive}
~
    \begin{exerpart}
        \item Construct a truth table for $P \land (Q \lor R)$.  Include an intermediate column for $Q \lor R$.
        \item Construct a truth table for $(P \land Q) \lor (P \land R)$. Include intermediate columns for $P \land Q$ and for $P \land R$.
        \item What do you notice? % Distributive property
    \end{exerpart} 
\end{exer}

\begin{exer}[\exerR] % And and Or
\label{exer_dyk_and_or}
    Do you know \ldots
        \begin{exerpart}
            \item What the symbols $\land$, $\lor$, and $\oplus$ stand for?
            \item What the logical connectives ``and'', ``or'', and ``xor'' mean? 
            \item When statements of the form $P \land Q$ and $P \lor Q$ are true and when they're false?
            \item Why we might create a truth table?
            \item How to create a truth table for a compound logical statement? 
        \end{exerpart}
\end{exer}

%%%%%%%%%%%%%%%%%%%%%%%%

\subsubsection*{Exercises for \Cref{sub_pff_and_or} \subpffandor}

\begin{exer}[\exerP] % And divides mult
\label{exer_prove_odd_and_mult3}
    Use \Cref{pff_and} to prove that 12 is even and $3 \mid 12$. Hint: Copy the proof in \Cref{act_prove_and_orusingcases}  (\ref{part_pff_and_example}) and edit as needed.   
\end{exer}

\begin{exer}[\exerP] % square+n mod 4
\label{exer_prove_nsqplusn_mod4}
    Use \Cref{pff_or_by_cases} to to prove that $n^2+2n \fn{ mod } 4 = 0$ or $n^2+2n \fn{ mod } 4 = 3$ for any integer $n$.  Hint: Copy the proof in \Cref{act_prove_and_orusingcases} (\ref{part_pff_or_using_cases_example}) and edit as needed.  Yes, this exercise repeats \Cref{act_prove_and_orusingcases} (\ref{part_nsqplusn_mod4}).
\end{exer} 

\begin{exer}[\exerU] % And divides mult general
\label{exer_prove_odd_and_mult3_general}
    Use \Cref{pff_and} to prove that $3n+6$ is even and $3 \mid (3n+6)$ for any integer $n$. Hint: Copy the proof in \Cref{act_prove_and_orusingcases}  (\ref{part_pff_and_example}) and edit as needed.    
\end{exer}

\begin{exer}[\exerU] % Prove squares mod 5
\label{exer_sq_mod5}
    Use Proof Format \Cref{pff_or_by_cases_generalized} to prove that for any integer $n$ we have $n^2 \fn{ mod } 5 = 0$, $n^2 \fn{ mod } 5 = 1$, or $n^2 \fn{ mod } 5 = 4$.  Hint: By the Division Algorithm $n = 5q+r$ with $0 \le r < 5$.  Break into five cases corresponding to $r=0, 1, 2, 3$, or $4$.
\end{exer}

\begin{exer}[\exerR] % Pff and and or by cases
\label{exer_dyk_pff_and_or_by_cases}
    Do you know \ldots
        \begin{exerpart}
            \item How we prove a statement of the form $P \land Q$?
            \item How we prove a statement of the form $P \lor Q$ using cases? 
        \end{exerpart}
\end{exer}

%%%%%%%%%%%%%%%%%%%%%%%%

\subsubsection*{Exercises for \Cref{sub_logical_equivalence_negations} 
\subequivneg}

\begin{exer}[\exerP] % negate statements
\label{exer_my_first_negate}
    Negate each statement about an integer $k$.
        \begin{exerpart}
            \item $k=10$
            \item $k>10$
            \item $k$ is odd
            \item $10 \mid k$          
        \end{exerpart}
\end{exer}

\begin{exer}[\exerP] % negate and/or statements
\label{exer_my_first_demorgans}
    Negate each statement about the integers $a$ and $b$.  Hint: Use DeMorgan's Laws (\Cref{thm_demorgans}) and \Cref{rem_but_not}.
        \begin{exerpart}
            \item $a$ is even and $b$ is even
            \item $a$ is even or $b$ is even  
            \item $a$ is even but $b$ is odd
        \end{exerpart}
\end{exer}

\begin{exer}[\exerU] % translating into statements using $\land$, $\lor$, $\oplus$ and $\neg$]
\label{exer_translate_into_and_or_with_not}
    For this exercise, $n$ is a specific (but unknown) integer,       
        \begin{center}
            $X=$ $n$ is a multiple of 6, and $R=$ $n$ is a multiple of 4.
        \end{center}  
    In each part of this exercise write the statement using $X$, $R$, and logic symbols and find an example of an integer $n$ for which the statement is true.
        \begin{exerpart}
            \item The integer $n$ is not a multiple of 6.
            \item The integer $n$ is a multiple of 6, but not a multiple of 4.
            \item The integer $n$ is a multiple of 4, but not a multiple of 6.   
            \item The integer $n$ is not a multiple of 4 or 6.               
        \end{exerpart}
\end{exer}

\begin{exer}[\exerP] % truth tables for ``and/or not'']
\label{exer_truthtable_andornot}
~
    \begin{exerpart}
        \item Use a truth table to show that $P \lor \neg P$ is a tautology. 
        \item Draw a truth table for $P \lor \neg Q$.  Include an intermediate column for $\neg Q$. 
        \item Draw a truth table for $\neg P \land Q$.  Include an intermediate column for $\neg P$.  
    \end{exerpart}
\end{exer}

\begin{exer}[\exerU] % verify DeMorgan's Laws]
\label{exer_truthtable_demorgan}
   In \Cref{act_not_and} we verified the first of De Morgan's Laws (\Cref{thm_demorgans}). Construct truth tables to verify the second equivalence $$\neg(P \lor Q) \equiv \neg P \land \neg Q.$$  Include some intermediate columns. 
\end{exer} 

\begin{exer}[\exerR] % Logical Equivalence and Negation
\label{exer_dyk_equiv_negation}
    Do you know \ldots
        \begin{exerpart}
            \item What it means for statements to be logically equivalent and how we might use a truth table to verify an equivalence?
            \item What a tautology or contradiction is and how to verify that a statement is a tautology or contradiction using a truth table. 
            \item What the negation of a statement is?
            \item What the symbols $\neg$ and $\equiv$ mean?
            \item How the construction $\land \neg$ is translated into words. 
            \item How to find the negation of $P \land Q$ and $P \lor Q$ and what those equivalences are called?
        \end{exerpart}
\end{exer}

\begin{exer}[\exerE] % writing $P \oplus Q$ without $\oplus$]
\label{exer_writing_xor_without_xor}
~   
    \begin{exerpart}
        \item Find a statement that is logically equivalent to $P \oplus Q$ but does not use $\oplus$. You may use $\land$, $\lor$, $\neg$, and parentheses. 
        \item Use a truth table to verify that the statements are equivalent.
        \item What is the negation of the statement $P \oplus Q$?  Describe the negation using $\land$, $\lor$, and $\neg$.
    \end{exerpart}
\end{exer}

%%%%%%%%%%%%%%%%%%%%%%%%%%

\subsubsection*{Exercises for \Cref{sub_pff_contra} \subpffcontra}

\begin{exer}[\exerP] % there do not exist three distinct digits whose sum is 25]
\label{exer_prove_not_sum_3digits_equals25}
    Use a proof by contradiction (\Cref{pff_contradiction}) to prove that there do not exist three distinct digits whose sum is equal to 25.  This statement was one of the true/false questions in \Cref{act_TorF} (\ref{part_digit_sum25}). Hint: The negation would be that there exist three distinct digits whose sum is 25.  What is the largest possible set of three distinct digits and what is their sum?
\end{exer}

\begin{exer}[\exerP] % Pidgeonhole principle
\label{exer_contradiction_coins_in_cups}
    We have nine coins and four cups. We place each coin in a cup. Use \Cref{pff_contradiction} to prove that some cup has three or more coins.  Hint: The negation would be that each cup has two or fewer coins.
\end{exer}

\begin{exer}[\exerP] % odd 3-regular graph
\label{exer_contradiction_3reg_graph_order7}
    Use a proof by contradiction (\Cref{pff_contradiction}) to prove that there cannot exist a 3-regular graph of order 7. Hint: By \Cref{thm_first_thm_graph_thy}, the number of edges in a graph is equal to half of the sum of the degrees of the vertices. On the other hand, the number of edges must be an integer.
\end{exer}

\begin{exer}[\exerU] % Every bit string of length five has at least three \bs{0}s or at least three \bs{1}s.]
\label{exer_prove_bs_length5_three0s_or_three1s}
~
    \begin{exerpart}
        \item Use a proof by contradiction (\Cref{pff_contradiction}) to prove that every bit string of length five has at least three \bs{0}s or at least three \bs{1}s.  Hint: The negation would be that there exists a bit string of length five with fewer than three $\bs{0}$s and fewer than three $\bs{1}$s. What is the longest such possible bit string?
        \item Why does the negation use ``and''?
    \end{exerpart}
\end{exer} 

\begin{exer}[\exerU] % coprime integers have no common prime divisor
\label{exer_contradiction_coprime}
     Suppose $a$ and $b$ are coprime integers and $p$ is a positive prime. Use a proof by contradiction (\Cref{pff_contradiction}) to prove that $p \nmid a$ or $p \nmid b$. Hint: Be careful negating the ``or'' statement.
\end{exer}

\begin{exer}[\exerR] % Pff Proof by Contradiction
\label{exer_dyk_pff_contradiction}
    Do you know \ldots
        \begin{exerpart}
            \item Why $Q \land \neg Q$ is a contradiction?
            \item How to write a proof by contradiction?
            \item How we can begin a proof by contradiction so the reader knows that we are doing a proof by contradiction?
            \item Why the phrase ``On the other hand'' is used in a proof by contradiction?
        \end{exerpart}
\end{exer}

\begin{exer}[\exerE] % cannot be both even and odd
\label{exer_not_both_even_odd}
    Use a proof by contradiction (\Cref{pff_contradiction}) to prove that an integer $n$ cannot be even and odd. Hint: If $n$ is even, use the definition of even to write $n = 2k$ for some integer $k$.  If $n$ is odd, use the definition of odd to write $n=2q+1$ for some integer $q$. Do not use $k$ again.  Put these equations together to get $2k=2q+1$.  Next, solve for 1 to show that $2 \mid 1$.  On the other hand, the only divisors of 1 are $\pm1$.
\end{exer}

\begin{exer}[\exerE] % lg(3) is irrational
\label{exer_lg3_irrat}
    Copy the following proof and fill in the blanks.\footnote{Thanks to Oscar Levin for the idea for this exercise.}
    
    Prove that $\fn{lg}(3)$ is irrational. 

    \begin{proof}
        Suppose not.  Then $\fn{lg}(3)$ is \underline{\hspace{.5 in}}.  By the definition of rational numbers (\Cref{defn_fraction_rational}), we can write 
            $$\fn{lg}(3) = \frac{s}{t}$$ 
        for some integers $s$ and $t$ with $t \neq 0$.  Using the definition of \fn{lg} (\Cref{defn_log}), we can write
            $$2^{^{\frac{s}{t}}} = \underline{\hspace{.5 in}} .$$
        Raising both sides of this equation to the power $t$ we get 
            $$2^s = \underline{\hspace{.5 in}} .$$

        On the other hand, any power of two is even and so \underline{\hspace{.5 in}} is even and any power of three is odd and so \underline{\hspace{.5 in}} is odd.  By \Cref{{exer_not_both_even_odd}}, an integer cannot be even and odd.  Thus, 
            $$2^s \neq \underline{\hspace{.5 in}}.$$  

        This is a contradiction.  Thus, $\fn{lg}(3)$ is irrational.
    \end{proof}
\end{exer}
\begin{exer}[\exerE] % prove there are infinitely many primes]
\label{exer_proof_inf_many_primes}
~
    Copy the following proof and fill in the blanks.
    
    Prove that there are infinitely many primes.

        \begin{proof}
            Suppose not.  Then there would be a \underline{\hspace{1in}} number of primes, so we could list all of them. Suppose that the list of all primes were $p_1,p_2,p_3, \ldots,p_k.$ Consider the positive integer $$n = p_1p_2p_3\cdots p_k.$$  

            First, $n = p_1(p_2p_3\cdots p_k)$ so by the definition of $\underline{\hspace{1in}}$ it follows that $p_1 \mid n$.  Since $p_1$ is prime, we know that $p_1 > 1$.  By \Cref{exer_consec_coprime}, we know that $\fn{gcd}(n,n+1) = \underline{\hspace{1in}}$. Since $p_1 > 1$ and $p_1 \mid n$, it follows that $p_1 \nmid (n+1)$.  For each value of $i = 1, 2, \ldots k$, the same argument shows that $p_i >1$ and $p_i \mid n$ and so $p_i \nmid (n+1)$. Therefore, none of the primes $p_1,p_2,p_3, \ldots, p_k$ divides $\underline{\hspace{1in}}$.  
            
            On the other hand, $p_1,p_2,p_3, \ldots,p_k$ are the only primes.  It follows that $n+1$ has no prime divisors. Since $n$ is positive, we know $n+1>1$. This is a contraction of \Cref{thm_FTOArith} which says $\underline{\hspace{1in}}$.
        \end{proof}
\end{exer}

%%%%%%%%%%%%%%%%%%%%%%%%%%%
\subsection{\answers}
%%%%%%%%%%%%%%%%%%%%%%%%%%%

%%%%%%%%%%%%%%%%%%
\subsubsection{\answers for \Cref{sub_and_or} \subandor}

\subsubsection*{\Cref{exer_translate_into_and_or} \answers}
\begin{apart}
    \item $X \land R$ which is true when $n=12$. There are other correct examples.
    \item Hint: $\oplus$
\end{apart}

\subsubsection*{\Cref{exer_truthtable_distributive} \answers}
\begin{apart}
    \item \label{part_truthtable_dist1} See \Cref{table_truthtable_dist1}.
            \begin{table}[hbtp]
                \centering
                \begin{tabular}{ccc |c|c}
                    $P$ & $Q$ & $R$ & $Q \lor R$ & $P \land (Q \lor R)$ \\ \hline
                    \str{T} & \str{T} & \str{T} & \str{T} & \str{T}\\
                    \str{T} & \str{T} & \str{F} & \str{T} & \str{T}\\
                    \str{T} & \str{F} & \str{T} & \str{T} & \str{T}\\
                    \str{T} & \str{F} & \str{F}& \str{F} & \str{F}\\
                    \str{F} & \str{T} & \str{T} & \str{T} & \str{F}\\
                    \str{F} & \str{T} & \str{F} & \str{T} & \str{F}\\
                    \str{F} & \str{F} & \str{T} & \str{T} & \str{F}\\
                    \str{F} & \str{F} & \str{F}& \str{F} & \str{F}
                \end{tabular}
                \caption{Truth table for the compound statement $P \land  (Q \lor R)$.}
                \label{table_truthtable_dist1}
            \end{table}
    \item Hint: You should get the same final column as in (\ref{part_truthtable_dist1}).
    \item They are logically equivalent.
\end{apart}

%%%%%%%%%%%%%%%%%%
\subsubsection{\answers for \Cref{sub_pff_and_or} \subpffandor}

\subsubsection*{\Cref{exer_prove_odd_and_mult3} \answers}
    Hint: $12 = 2(6)$ and $12 = 3(4)$.

\subsubsection*{\Cref{exer_prove_nsqplusn_mod4} \answers}
    Hint: $$(2k)^2+2(2k) = 4k^2+4k = 4(\underline{\hspace{.5in}}) + 0$$ and $$(2k+1)^2 + 2(2k+1) = 4k^2 + 8k + 3= 4(\underline{\hspace{.5in}}) + 3.$$

\subsubsection*{\Cref{exer_prove_odd_and_mult3_general} \answers}
    Hint: Make sure to copy all the words in the proof.  During the proof, you will need to write $3n+6 = 2(\underline{\hspace{.5in}})$ and $3n+6 = 3(\underline{\hspace{.5in}}).$

\subsubsection*{\Cref{exer_sq_mod5} \answers}
    Hint: The five cases are Case 1: Assume $n = 5q$; Case 2: Assume $n=5q+1$; Case 3: Assume $n=5q+2$; Case 4: Assume $n=5q+3$; and Case 5: Assume $n=5q+4$. In Case 4, for example, we have 
        $$n^2=(5q+3)^2=25q^2+30q+9=25q^2+30q+5+4=5(5q^2+6q+1)+4$$ 
    and so $n^2 \fn{ mod } 5 = 4$.  Notice how we cleverly rewrite $9 = 5+4$.  That made it possible to factor 5 out of the 5 to show the remainder of 4.

%%%%%%%%%%%%%%%%%%
\subsubsection{\answers for \Cref{sub_logical_equivalence_negations} 
\subequivneg}

\subsubsection*{\Cref{exer_my_first_negate} \answers}
\begin{apart}
    \item $k \neq 10$
    \item Hint: Make sure your answer includes the possibility that $k=10$.
    \item $k$ is even.
    \item Hint: Use $\nmid$.
\end{apart}

\subsubsection*{\Cref{exer_my_first_demorgans} \answers}
\begin{apart}
    \item $a$ is odd or $b$ is odd
    \item Hint: Your answer should include ``and''.
    \item Hint: Your answer should include ``or''.
\end{apart}

\subsubsection*{\Cref{exer_translate_into_and_or_with_not} \answers}
\begin{apart}
    \item $\neg X$  An example is $n=10$.  There are other correct examples.
    \item Hint: See \Cref{rem_but_not} and use $\nmid$. Include an example.
    \item Hint: $R \land \neg X$.  Include an example.
    \item Hint: You may use parentheses in your answer, but there is also a way to write it without parentheses. Include an example.
\end{apart}

\subsubsection*{\Cref{exer_truthtable_andornot} \answers}
\begin{apart}
    \item Hint: The truth table only needs two rows and is similar to the table in \Cref{table_truthtable_contradiction}.
    \item Hint: The truth table needs four rows and is similar to the table in \Cref{table_truthtable_andnot}. The last column should have three \str{T} and one \str{F}.
    \item Hint: The last column should have one \str{T} and three \str{F}.
\end{apart}

\subsubsection*{\Cref{exer_truthtable_demorgan} \answers}
  Hint: One way to set up the truth table is shown in \Cref{table_truthtable_demorgan}. Feel free to copy the table, check the third row, and then fill in the remaining entries.
            \begin{table}[hbtp]
                \centering
                \begin{tabular}{cc ||c|c||c|c|c}
                    $P$ & $Q$ & $P \lor Q$ & $\neg(P \lor Q)$ & $\neg P$ & $\neg Q$ & $\neg P \land \neg Q$\\ \hline
                    \str{T} & \str{T} &&&&& \\ \hline
                    \str{T} & \str{F} &&&&& \\ \hline
                    \str{F} & \str{T} & \str{T} & \str{F} &\str{T} & \str{F} & \str{F} \\ \hline
                    \str{F} & \str{F} &&&&& \\
                \end{tabular}
                \caption{Truth table for the compound statement $P \land  (Q \lor R)$.}
                \label{table_truthtable_demorgan}
            \end{table}

%%%%%%%%%%%%%%%%%%%%%%%%%%
\subsubsection{\answers for \Cref{sub_pff_contra} \subpffcontra}

\subsubsection*{\Cref{exer_prove_not_sum_3digits_equals25} \answers}
    Hint: Copy this proof filling in the missing parts.
    \begin{proof}
        Suppose not.  Then there exist three distinct digits whose sum is equal to 25.

        On the other hand, the largest possible set of three distinct digits is $\{ \underline{\hspace{.1in}}, \underline{\hspace{.1in}}, \underline{\hspace{.1in}}\}$  and their sum is $\underline{\hspace{.1in}} + \underline{\hspace{.1in}}+\underline{\hspace{.1in}}=\underline{\hspace{.1in}} \neq 25$.

        This is a contradiction.  Therefore, there do not exist three distinct digits whose sum is equal to 25.
    \end{proof}

\subsubsection*{\Cref{exer_contradiction_3reg_graph_order7} \answers}
    Hint: If there were such a graph, the sum of the degrees of the vertices would be     
        $$3+3+3+3+3+3+3=7(3)=21.$$
    Then, use the hint.

\subsubsection*{\Cref{exer_prove_bs_length5_three0s_or_three1s} \answers}
\begin{apart}
    \item Hint: In the hint ``fewer than three \bs{0}s'' means the number of \bs{0}s is less than or equal to two.
    \item Hint: DeMorgan's Law.
\end{apart}

\subsubsection*{\Cref{exer_contradiction_coprime} \answers}
  Hint: If $p \mid a$ and $p \mid b$, then $p$ is a positive common divisor of $a$ and $b$ that is greater than one.  What does that tell you about $\fn{gcd}(a,b)$?





